% !TeX spellcheck = es_ES
\documentclass[letterpaper,11pt]{article}
\usepackage[utf8]{inputenc}
\usepackage[T1]{fontenc}
\usepackage[spanish]{babel}
\decimalpoint % separador decimal con punto, como en las tablas de la fuente original
\usepackage{lmodern,microtype}
\usepackage{amsmath,amssymb,amsfonts}
\usepackage{geometry}
\geometry{margin=2.54cm}
\usepackage{setspace}
\setstretch{1.15}
\usepackage{enumitem}
\usepackage{booktabs}
\usepackage{color}
\usepackage{hyperref}
\definecolor{ntr}{rgb}{0.8,0,0}
\hypersetup{colorlinks=true,linkcolor=ntr,citecolor=ntr,urlcolor=black}

% Portada (Luis Cabral): titulo \LARGE sans-serif alineado a la izquierda, sin numero de pagina
\makeatletter
    \renewcommand{\@maketitle}{%
    \newpage\parindent0in\null\vskip2em%
    {\LARGE\textbf{\textsf{\@title}}\par\vskip2em}%
    \@author\par\vskip2em%
    \@date%
    \par
    \thispagestyle{empty}\setcounter{page}{1}%
    }%
\makeatother

\DeclareMathOperator{\Var}{Var}
\DeclareMathOperator{\Cov}{Cov}
\DeclareMathOperator{\Corr}{Corr}
\DeclareMathOperator{\E}{\mathbb{E}}
\DeclareMathOperator{\N}{Normal}
\DeclareMathOperator{\LP}{\mathbb{L}}
\newcommand\independent{\protect\mathpalette{\protect\independenT}{\perp}}
\def\independenT#1#2{\mathrel{\rlap{$#1#2$}\mkern2mu{#1#2}}}
% Teoria asintotica
\newcommand{\SR}{\mathbb{R}}
\DeclareMathOperator{\plim}{plim}
\DeclareMathOperator{\Avar}{AVar}
\newcommand{\pto}{\stackrel{p}{\longrightarrow}}
\newcommand{\dto}{\stackrel{d}{\longrightarrow}}
\newcommand{\adistr}{\stackrel{a}{\sim}}
\newcommand{\epsi}{\varepsilon}


\title{Análisis Econométrico \\ Ayudantía 11}
\author{Magíster en Economía, Universidad Alberto Hurtado}
\date{}

% Funciones de R sugeridas: entre corchetes, alineadas a la derecha; si no caben en la
% linea pasan enteras a la siguiente.
\newcommand{\fn}[1]{{\unskip\nobreak\hfill\penalty50\hskip1em\null\nobreak\hfill
    \mbox{[#1]}\parfillskip=0pt\finalhyphendemerits=0\par}}

\begin{document}
\maketitle

\noindent Temas: supuestos de variables instrumentales (relevancia, exogeneidad y
no colinealidad) y qué pasa cuando no se cumplen; condición de rango y relevancia conjunta con
varios instrumentos; validez de un instrumento en aplicaciones; y los contrastes de
sobreidentificación (Sargan) y de endogeneidad (Wu--Hausman) calculados a mano en R.

\medskip
\noindent Los ejercicios 1 a 3 se resuelven en la pizarra. El ejercicio 4 es una demostración en
R: el ayudante corre el código en pantalla, así que no hace falta traer computador.

\section*{Si quiere replicar el ejercicio 4 en su computador}

\noindent El script \texttt{TA11.R} y los datos están en el repositorio del curso, en
\href{https://github.com/ramirode/econometria_mae/tree/main/TA/TA11}{\texttt{TA/TA11}}. Los pasos
son los mismos de la Ayudantía 5: cree una carpeta \texttt{TA11} con las subcarpetas
\texttt{code} (para \texttt{TA11.R}) e \texttt{input} (para \texttt{bancos.rds}, los mismos datos
de la Ayudantía 10), y cambie la línea \texttt{dir\_proyecto <- "..."} por la ruta a su carpeta
\texttt{TA11}. Paquetes: \texttt{sandwich}, \texttt{lmtest}, \texttt{car} e \texttt{ivreg}.

\section*{Ejercicios}

\begin{enumerate}[itemsep=0.8em]

%% ---------------------------------------------------------------- 1
\item Dado un modelo con una variable explicativa exógena $x_1$ y una variable explicativa
    endógena $x_2$,
    \begin{align*}
    y=\beta_0+\beta_1 x_1+\beta_2 x_2+u,
    \end{align*}
    donde $\E(u)=0$, $\E(x_1u)=0$ pero $\E(x_2u)\neq0$. Suponga que tiene un instrumento $z_1$
    disponible. Sea $z=(1,x_1,z_1)'$ el vector de variables exógenas y $x=(1,x_1,x_2)'$. Los
    supuestos son: (i) exogeneidad, $\E(zu)=0$; (ii) relevancia,
    $\operatorname{rango}\E(zx')=3$; y (iii) no multicolinealidad perfecta entre las exógenas,
    $\operatorname{rango}\E(zz')=3$. Explique qué supuesto \emph{no} se cumple en cada uno de los
    siguientes casos (y por qué):
    \begin{enumerate}
        \item $z_1$ es independiente de $(y,x_1,x_2)$.
        \item $z_1=x_2$.
        \item $z_1=1$.
        \item $z_1=x_1$.
    \end{enumerate}

%% ---------------------------------------------------------------- 2
\item Considere el modelo del ejercicio 1, pero ahora con dos instrumentos disponibles,
    $z_1$ y $z_2$, y sea $z=(1,x_1,z_1,z_2)'$ el vector de variables exógenas.
    \begin{enumerate}
        \item Escriba la primera etapa de $x_2$,
            $x_2=\pi_{0}+\pi_{1}x_1+\pi_{2}z_1+\pi_{3}z_2+e$. Muestre que la condición de rango,
            $\operatorname{rango}\E(zx')=3$ con $x=(1,x_1,x_2)'$, se cumple si y sólo si
            $\pi_2\neq0$ o $\pi_3\neq0$. ¿Cómo contrastaría la condición de rango?
        \item Suponga que $\pi_2\neq0$ pero $\pi_3=0$. ¿Está identificado el modelo? ¿Qué se
            pierde, o se gana, al incluir $z_2$ como instrumento?
        \item Suponga ahora que $x_1$ \emph{también} es endógena y que se usan los mismos dos
            instrumentos $z_1$ y $z_2$ para las dos variables endógenas; en este caso
            $z=(1,z_1,z_2)'$. Si en la primera etapa de $x_1$ y en la de $x_2$ los dos
            instrumentos son individual y conjuntamente significativos, ¿puede estar seguro de
            que el modelo está identificado? Dé un ejemplo en el que no lo esté.
    \end{enumerate}

%% ---------------------------------------------------------------- 3
\item Para cada uno de los siguientes estudios, (i) explique por qué MCO puede ser inconsistente
    y sugiera el signo del sesgo, (ii) explique por qué el instrumento puede ser relevante y
    (iii) describa un canal por el cual el instrumento podría no ser exógeno y qué le haría esto
    a la estimación por variables instrumentales.
    \begin{enumerate}
        \item \emph{Servicio militar e ingresos} (Angrist, \textit{American Economic Review}, 1990).
            Efecto de haber servido en el ejército durante la guerra de Vietnam sobre los
            ingresos posteriores. A comienzos de los años setenta EE.UU. asignó por sorteo, según la fecha de
            nacimiento, un número a cada joven; los jóvenes con número bajo
            eran elegibles para ser reclutados y los demás no. Instrumento: elegibilidad según el
            número del sorteo.
        \item \emph{Instituciones y desarrollo} (Acemoglu, Johnson y Robinson,
            \textit{American Economic Review}, 2001). Efecto de la calidad de las instituciones
            (protección contra la expropiación) sobre el ingreso per cápita de los países hoy.
            Instrumento: la mortalidad de los colonos europeos en cada colonia en el siglo XIX.
        \item \emph{Policía y crimen} (Levitt, \textit{American Economic Review}, 1997). Efecto
            del número de policías por habitante sobre la tasa de crimen violento en las ciudades
            de EE.UU. Instrumento: si hubo elecciones de alcalde o gobernador ese año, porque los
            políticos contratan más policías en años electorales.
    \end{enumerate}

%% ---------------------------------------------------------------- 4
\item \emph{Sargan y Wu--Hausman a mano.} Aún no hemos visto estos dos contrastes en clase,
    así que primero una introducción breve. En ambos, $z$ es el vector de todas las variables
    exógenas (los controles y los instrumentos), $L$ es su número y $K$ el número de coeficientes
    del modelo.
    \begin{itemize}
        \item \emph{Sargan} contrasta la exogeneidad de los instrumentos cuando hay más
            instrumentos que endógenas, es decir, $L-K>0$. $H_0$: $\E(zu)=0$. Si los instrumentos
            son exógenos, los residuos de 2SLS, $\hat u_i$, no deberían poder predecirse con $z$.
            Se regresa $\hat u_i$ en $z$ y se calcula $S=N\cdot R^2$, que bajo $H_0$ (y con
            errores homoscedásticos) se distribuye $\chi^2_{L-K}$. Si $L=K$, $R^2=0$ por
            construcción y no hay nada que contrastar. Si rechaza, al menos uno de los
            instrumentos no es exógeno. Si no rechaza, no queda demostrado que lo sean.
        \item \emph{Wu--Hausman} contrasta si una variable explicativa es endógena.
            $H_0$: $\E(x_ju)=0$. Si $H_0$ es cierta, MCO y 2SLS son consistentes y deberían dar
            estimaciones parecidas; si difieren, $x_j$ es endógena. Una forma de calcularlo
            (función de control): (i) se estima la primera etapa de $x_j$ y se guardan los
            residuos $\hat v_i$; (ii) se estima por MCO la ecuación de interés agregando
            $\hat v_i$ como regresor; (iii) se contrasta que el coeficiente de $\hat v_i$ es
            cero con un test $t$. Si rechaza, hay endogeneidad y conviene usar el estimador de
            variables instrumentales. Si no rechaza, MCO y VI no son distinguibles, pero eso
            tampoco prueba que $x_j$ sea exógena.
    \end{itemize}
    Con los datos \texttt{bancos.rds} del ejercicio 2 de la
    Ayudantía 10 (Pascali, 2016), considere el modelo
    \begin{align*}
    pib_{it}=\beta_0+\beta_1\, credito_{it}+x_{it}'\gamma+u_{it},
    \end{align*}
    donde $x$ incluye las características geográficas del municipio y dummies de provincia y de
    año, e instrumente $credito$ con $jew\_peq$, $jew\_med$ y $jew\_gde$.
    \begin{enumerate}
        \item Estime por 2SLS con \texttt{ivreg} y guarde los tests de Sargan y de Wu--Hausman que
            entrega \texttt{summary(\dots, diagnostics = TRUE)}.
            \fn{\texttt{ivreg}, \texttt{summary}}
        \item \emph{Sargan a mano.} Guarde los residuos de 2SLS, $\hat u_i$. Regréselos en todas
            las exógenas (los controles y los tres instrumentos). ¿Cuántas restricciones de
            sobreidentificación hay? Calcule $N\cdot R^2$ y compárelo con el valor crítico de la
            $\chi^2$ con ese número de grados de libertad.
            \fn{\texttt{resid}, \texttt{lm}, \texttt{pchisq}}
        \item \emph{Wu--Hausman a mano.} Estime la primera etapa de $credito$ (con los tres
            instrumentos) y guarde los residuos $\hat v_i$. Estime por MCO la ecuación de interés
            agregando $\hat v_i$ como regresor. Contraste que el coeficiente de $\hat v_i$ es
            cero con errores estándar robustos. Compare el estadístico $t^2$ con el de (a), y el
            coeficiente de $credito$ con el de 2SLS. ¿Qué dice el signo del coeficiente de
            $\hat v_i$?
            \fn{\texttt{lm}, \texttt{vcovHC}, \texttt{coeftest}}
        \item El test de Sargan no rechaza. ¿Significa que los tres instrumentos son
            exógenos? ¿Se podría hacer este test con un solo instrumento?
    \end{enumerate}

\end{enumerate}

\end{document}
